\scanpage{152}
\mspage{75}
\uline{Corollaire 2} On a
\[
(246)\qquad \Pi_1\underbrace{\deleted{X}^{*\infty}}_{=V^*_{\Theta,X}}
\simeq\underset{\longrightarrow}{\operatorname{Lim}}
\left(\vphantom{\widetilde\Pi_{2,1}}\right.
\]
\diagram{\begin{tikzcd}[ampersand replacement=\&,column sep=small]
\&\widetilde\Pi_{2,1}\arrow[dl]\arrow[dr]\&\\
\widetilde\Pi_1\&\&\widetilde\Pi_2
\end{tikzcd}}\[\left.\vphantom{\widetilde\Pi_{2,1}}\right).\]

En effet, une analyse locale facile (du type ``th. de pureté'') montre qu'on peut innocemment négliger $\Theta^{(3)}$, i.e. que ça \unsure{pour nous} ne change pas si l'on remplace $X$ par $X\setminus\Theta^{(3)}$.

\margin{Il faut expliciter aussi comment envoyer les $\widetilde\Pi_i$ $(i\geq3)$ dans le groupoïde $\underset{\longrightarrow}{\operatorname{Lim}}$, p.ex. $\widetilde\Pi_3$ est bien $\simeq$ \unsure{représenté|reconstitué} comme groupoïde $\underset{\longrightarrow}{\operatorname{Lim}}$ du
\diagram{\begin{tikzcd}[ampersand replacement=\&,column sep=small]
\&\widetilde\Pi_{321}\arrow[dl]\arrow[dr]\&\\
\widetilde\Pi_{3,1}\&\&\widetilde\Pi_{3,2}
\end{tikzcd}}, formé de $2$ groupoïdes qui s'envoient dans
$\underset{\longrightarrow}{\operatorname{Lim}}$
\diagram{\begin{tikzcd}[ampersand replacement=\&,column sep=small]
\&\widetilde\Pi_{2,1}\arrow[dl]\arrow[dr]\&\\
\widetilde\Pi_1\&\&\widetilde\Pi_2
\end{tikzcd}}.
Il faudra expliciter les \unsure{2-flèches}
\diagram{\begin{tikzcd}[ampersand replacement=\&,column sep=tiny]
\&\unsure{\cdot}\arrow[dl]\arrow[dr]\&\\
\unsure{\cdot}\arrow[rr,dotted]\&\&\unsure{\cdot}
\end{tikzcd}}.}

\uline{Corollaire 3} Supposons $X^{*\infty}\longrightarrow X^*$ soit une $1$-équivalence, i.e induise une équivalence pour les $\Pi_1$. Alors
\[
(247)\qquad \Pi_1X\simeq\underset{\longrightarrow}{\operatorname{Lim}}
\]
\diagram{\begin{tikzcd}[ampersand replacement=\&,column sep=small]
\&\widetilde\Pi_{\unsure{2,1}}\arrow[dl,"\uline{générisation}"']\arrow[dr]\&\\
\widetilde\Pi_1\&\&\widetilde\Pi_2
\end{tikzcd}}.

\deleted{\uline{Corollaire 4} \inserted{\uline{\unsure{Sans déf. N, etc.}}} Supposons que pour \deleted{\unsure{tt|tout}} $i\leq d$, l'on ait
\diagram{\begin{tikzcd}[ampersand replacement=\&]
D_i^{*\infty}\arrow[r,leftrightarrow]\arrow[d,equal,"\mathrm{def}"']\&1\\
\deleted{\unsure{D_i^*}}V^*_{\Theta_i,D_i}\setminus\Theta_i\&
\end{tikzcd}},
considérons une \deleted{partie} sous-multiplicité \inserted{$D'_i$} d'un $D_i$ qui soit ouverte et fermée, et considérons l'\unsure{image} $\Theta'^{(i)}\subset\Theta^{(i)}$ de $D'_i$, qui est donc une ``str. de codim $i$'',}
\margin{151}
